\section{Source Witnessing, Theorem Fate, and Publication Boundary}
\label{sec:oc13-publication-boundary}

This chapter fixes the publication boundary of Core~1.3. Its task is to state
which scientific object is canonical, how source witnessing is enforced, how
theorem fate is separated across outward artifacts, and how later channels must
derive from the flagship monograph. It does not re-argue the formal doctrine or
theorem spine; those burdens belong to the surrounding chapters and appendices.

\subsection{Canonical scientific object}

The canonical scientific object of the release is the present master
monograph. It is the volume in which the orientation ladder, formal doctrine,
theorem route, empirical discipline, synthesis, practical utility, and audit
layers remain visible together. Narrower outward artifacts may be derived from
it, but they are not allowed to displace it as the authority on claim scope,
notation, or proof boundary.

This rule matters because foundational programmes accumulate several different
textual layers over time: source-native definitions and theorems, explanatory
prose, repair notes, channel-specific summaries, and later operational packets.
Core~1.3 treats those layers as distinguishable, not interchangeable.

\subsection{Source witnessing}

Core~1.3 is source-first in the strict sense that the flagship claims used in
public release must remain bound to public source witnesses rather than to late
summary prose alone. The point is not archival fetishism. The point is to stop
drift between five things that often blur inside large research programmes:

\begin{itemize}
\item theorem-bearing statements,
\item structural and interpretive explanation,
\item repair and release notes,
\item condensed public packets,
\item and later operational projections.
\end{itemize}

When that boundary blurs, a summary sentence can begin to masquerade as the
theorem itself, or a release dashboard can begin to replace the actual
manuscript. Core~1.3 blocks that failure by keeping source witnessing inside
the flagship scientific object rather than relegating it to hidden control
files.

\subsection{Theorem fate and scope discipline}

The present release also separates theorem fate explicitly. Not every inherited
row serves the same role in every outward artifact.

\begin{itemize}
\item Some rows remain part of the formal backbone and survive unchanged.
\item Some remain scientifically useful but only under explicit scope notes.
\item Some belong in support layers, appendices, or narrower companion
      artifacts rather than in the positive body of every outward paper.
\item Some remain frontier or hypothesis material and therefore may not be
      silently promoted by summary prose.
\end{itemize}

This is not retreat. It is a condition of honest publication. The monograph
can remain large without rhetorical inflation only if the reader can tell which
rows are theorem-bearing, which are bounded, and which are support or frontier
material.

\subsection{Relation to journal-core and companion artifacts}

Core~1.3 therefore adopts a one-way publication logic:
\[
\text{Master Monograph} \longrightarrow \text{Derived outward artifacts}.
\]

The journal-core article is narrower by design. Companion packets, matrices,
and release surfaces are narrower or more operational by design. Those
artifacts may compress, subset, or reorganize the flagship material for a
specific task, but they may not enlarge the scientific claim, widen notation,
or outrun the theorem-fate discipline fixed here.

The practical consequence is simple. If an outward artifact says something
stronger than the flagship volume, the outward artifact is wrong, even if it is
shorter, cleaner, or easier to market.

\subsection{Downstream projection rule}

The same boundary stabilizes later use in review, collaboration, product,
business, and public communication. Different channels may legitimately cite
different parts of the flagship volume:

\begin{itemize}
\item the theorem route when bounded proof is the issue,
\item the empirical route when promotion discipline is the issue,
\item the practical route when use and comparison are the issue,
\item the audit appendices when challenge and traceability are the issue.
\end{itemize}

What they may not do is invent a second scientific object. Source witnessing,
theorem fate, and publication boundary exist so that external use can remain
concise without becoming interpretively lawless.

\paragraph{Bridge to the next chapter.}
With the publication boundary fixed, the next task is internal consistency:
which statement classes the flagship volume uses, why the hierarchy runs
through \(K_{12}\), and how upper levels remain lawful without rhetorical
inflation.

\clearpage
